\section{Local regularity of the minimum}\label{sec:localregularity}

We continue to regard the minimum as a curve in the complete all-norm
space. Its first regularity properties are metric and spectral; finite
apartments will be obtained later.

\subsection{Time allocation and strict velocity margins}
The barrier alone bounds speeds by $\pi/2$. To make local variations,
we need a strict bound, independent of the number of polygonal segments.

\begin{lemma}\label{lem:timeallocation}
Fix a compact interval $I=[l,r]$, its length $H>0$, and an ordered
polygon of action at most $E_0$. Its time allocation can be changed,
keeping the ordered norm nodes and endpoints, so that its action does
not increase and every spectral velocity has absolute value at most
$b<\pi/2$. The constant depends only on $I,E_0,n$, and $q$.
\end{lemma}

\begin{proof}
Minimize the exact action over the compact simplex of ordered node
times. A zero-duration segment has infinite cost unless its endpoints
coincide; in that case its cost is zero. Every nonzero edge has a
strictly positive minimum feasible duration and its cost diverges
at the barrier. Constant segments have continuous potential costs
as their duration vanishes. Thus the extended objective is lower
semicontinuous and attains a finite minimum. Delete zero-duration
segments. All remaining durations are positive, so the derivatives
with respect to interior node times vanish.

For segment $j$ put $\overline R_j=W_j/\Delta_j$, and let $v_{ji}$
range over all $2n$ spectral velocities. Define
\[
 L_j=\sum_i\ell(v_{ji}),\qquad H_j=\sum_i v_{ji}\tan v_{ji},
 \qquad Q_j=\tfrac12\overline R_j(H_j-L_j).
\]
If $g_j$ is the node determinant gap, then
$\kappa_j=(g_{j+1}-g_j)/\Delta_j$ is the difference of the two
velocity sums, hence $|\kappa_j|\le n\pi$. Set
\[
 b_{Lj}=\int_{u_j}^{u_{j+1}}R(u)\frac{u_{j+1}-u}{\Delta_j}\,\dd u,
 \qquad
 b_{Rj}=\int_{u_j}^{u_{j+1}}R(u)\frac{u-u_j}{\Delta_j}\,\dd u.
\]
Differentiating the exact kinetic and potential integrals at an
interior node, whose norm is held fixed, gives
\begin{equation}\label{eq:timeallocationidentity}
 \begin{split}
 Q_j-Q_{j-1}={}&\tfrac12(R(u_j)-\overline R_j)L_j
       -\tfrac12(R(u_j)-\overline R_{j-1})L_{j-1}\\
 &+q(\kappa_{j-1}b_{R,j-1}+\kappa_jb_{Lj}).
 \end{split}
\end{equation}
For example, the right-end kinetic derivative is
$R(u_{j+1})L_j/2-\overline R_jH_j/2$, and the right-end potential
derivative is $qR(u_{j+1})g_{j+1}-q\kappa_jb_{Rj}$.
The corresponding left-end derivatives cancel the endpoint potential
terms and give \eqref{eq:timeallocationidentity}.

Let $M=\max_I R$ and $S=\max_I|R'|$. Since the potential is
nonnegative and $R\ge1$, $\sum_j\Delta_jL_j\le2E_0$.
Using $|R(u_j)-\overline R_j|\le S\Delta_j$ and the analogous
bound on the preceding segment gives
\begin{equation}\label{eq:timeallocationvariation}
 \sum_j|Q_j-Q_{j-1}|\le2SE_0+2qn\pi MH.
\end{equation}
Some segment has $L_{j_0}\le 2E_0/H$. Every velocity on it has
absolute value at most $b_0=\arccos(e^{-2E_0/H})$, so
$Q_{j_0}\le Mn b_0\tan b_0$. Thus all $Q_j$ have a bound independent
of the number of segments. The function
\[
 f(v)=v\tan v-\ell(v)
\]
is even, nonnegative and strictly increasing for $v>0$, and tends
to infinity as $v\uparrow\pi/2$. Since
$Q_j\ge\tfrac12 f(v_{ji})$, the asserted strict bound follows.
\end{proof}

\begin{proposition}\label{prop:localmargin}
On every compact interval the vertex curves of the whole-line minimum
are $b$-Lipschitz in $d_\infty$ for some $b<\pi/2$. There is a global
minimizing sequence with polygonal restriction to that interval and
the same strict spectral velocity bound there.
\end{proposition}

\begin{proof}
Start with split polygonal almost-minima on expanding cutoff intervals,
attached to their reference tails. Proposition~\ref{prop:finitecutoff}
and regular-curve recovery make their $J$ values tend to $J_*$, so
they converge uniformly to $c_*$. For a fixed symmetric compact
interval $I=[-U,U]$, the truncated identity
\eqref{eq:renormalization} bounds their local $E_I$ values: their
truncated nonnegative $J_I$ integrals are bounded by $J$, and their
endpoint determinant gaps converge. Apply
Lemma~\ref{lem:timeallocation} only on $I$, leaving its endpoint
norms and the exterior unchanged. The change in $J$ is the change
in $E_I$, since all boundary terms cancel. The new curves are still
in $\calQ$ and still minimize asymptotically. They therefore have
the same uniform limit by \eqref{eq:globalstability}. Their uniform
strict bound on $I$ passes to that limit. An arbitrary compact
interval lies inside such an $I$.
\end{proof}

\subsection{Exact-endpoint local minimality}
The endpoints of a compact restriction of $c_*$ might still be nonsplit.
The following repair allows local competitors with those exact endpoints.

\begin{lemma}\label{lem:endpointrepair}
Suppose ordered polygons $p_m$ on $I$ converge uniformly to $c$ and
have a common strict spectral velocity margin. They can be replaced
by ordered all-norm polygons $\widehat p_m$ with the exact endpoints
of $c$, such that $\widehat p_m\to c$ uniformly and
$E_I(\widehat p_m)-E_I(p_m)\to0$.
Conversely, every fixed finite-action polygon $p$ with the endpoints
of $c$ has polygonal replacements $z_m$ with the endpoints of $p_m$
and $E_I(z_m)\to E_I(p)$.
\end{lemma}

\begin{proof}
Let $\delta_m$ be the maximum endpoint $d_2$ error. Choose
$\eta_m\downarrow0$ with $\delta_m/\eta_m\to0$ and
$2\eta_m<|I|$. Add product-geodesic bridges of length $\eta_m$
at the two ends, and compress $p_m$ affinely into the middle interval.
Ordered endpoints give ordered bridges. The middle velocities are
multiplied by $1+O(\eta_m)$ and retain a fixed strict margin; the
bridge velocities are bounded by $\delta_m/\eta_m$.

Change variables on the middle interval. Its Jacobian, time weight
and velocity scaling change by $O(\eta_m)$, while $\ell$ and $\ell'$
are uniformly bounded on the strict margin. Thus its kinetic change
is $O(\eta_m)$, independently of the number of segments. Uniformly
bounded determinant gaps give the same potential bound. Each bridge
has kinetic cost $O(\delta_m^2/\eta_m)=o(1)$ and potential cost
$O(\eta_m)$. The uniform Lipschitz bound and small time displacement
prove uniform convergence. Reverse the bridges and compress the
fixed polygon $p$ to prove the converse assertion. Its own finite
partition supplies the strict margin needed for the same estimates.
\end{proof}

\begin{proposition}\label{prop:localminimality}
Every compact restriction of $c_*$ is the unique exact-endpoint local
relaxed minimum of \eqref{eq:localrelaxation}.
\end{proposition}

\begin{proof}
Use the sequence from Proposition~\ref{prop:localmargin} on $I$.
The first endpoint repair gives
$E_I^*\le\liminf_m E_I(p_m|I)$. The local problem is feasible,
since the strict metric speed bound makes the endpoint chord have
finite cost. Take any local polygon $p$ with the exact limiting
endpoints. The converse repair gives $z_m$ with the split endpoints
of $p_m|I$ and $E_I(z_m)\to E_I(p)$.

The intermediate nodes of $z_m$ need not be split. Replace it by the
minimum on its finite partition. Proposition~\ref{prop:finitepartition}
makes this a split polygon $\widetilde z_m$ with the same endpoints
and no larger action. It can now be inserted into $p_m$ while retaining
membership in $\calQ$. Minimality of the whole-line infimum and
cancellation of the finite boundary terms give
\[
 E_I(p_m|I)\le E_I(\widetilde z_m)+J(p_m)-J_*
             \le E_I(z_m)+J(p_m)-J_*.
\]
Taking the upper limit and then the infimum over $p$ proves
$E_I(p_m|I)\to E_I^*$. Their exact-endpoint repairs have the same
limiting energy and converge uniformly to $c_*|I$.
Proposition~\ref{prop:localrelaxedminimum} identifies this curve
with the unique local minimum.
\end{proof}

\subsection{An ordered kinetic comparison}
Work over $L$, where all endpoint norms are split. Fix $I=[l,r]$ and
ordered endpoint pairs with
$d_\infty(h_v(l),h_v(r))<\pi|I|/2$. At each vertex choose a common
endpoint basis, with weights $a_{vi},b_{vi}$. There is a unique real
$p_{vi}$ satisfying
\begin{equation}\label{eq:bridgeparameter}
 \int_l^r\arctan(p_{vi}\sech u)\,\dd u=b_{vi}-a_{vi}.
\end{equation}
The left side is strictly increasing from $-\pi|I|/2$ to
$\pi|I|/2$. Define a split curve $k$ in these endpoint apartments by
\begin{equation}\label{eq:kineticbridge}
 w_{vi}(l)=a_{vi},\qquad
 w_{vi}'(u)=\arctan(p_{vi}\sech u),\qquad
 R(u)\tan w_{vi}'(u)=p_{vi}.
\end{equation}

\begin{lemma}\label{lem:kineticbridge}
The curve $k$ minimizes the pure kinetic action uniquely in its uniform
polygonal relaxation with these endpoints. It is ordered. If $\gamma$
is the affine endpoint chord, then
\begin{equation}\label{eq:bridgechord}
 \sup_{u\in I}d_2(k_v(u),\gamma_v(u))
                   \le\frac{\sqrt n}{16}|I|^2.
\end{equation}
\end{lemma}

\begin{proof}
Project a comparison polygon separately at the two vertices onto
their respective endpoint apartments using
Lemma~\ref{lem:blockprojection}. This fixes the endpoints and
decreases the kinetic cost. In either apartment, the supporting
inequality for $\ell$ and the constant momenta in
\eqref{eq:kineticbridge} give
\[
 \int_I\frac R2\sum_i\ell(z_i')\,\dd u
 \ge\int_I\frac R2\sum_i\ell(w_i')\,\dd u
       +\frac12\sum_i p_i\int_I(z_i'-w_i')\,\dd u.
\]
The final term vanishes by the endpoints. This proves the lower
bound against all polygons and, by recovery, identifies the kinetic
infimum with $K(k)$.

The midpoint calculation without potential gives, for two
exact-endpoint polygons $z,z'$,
\begin{equation}\label{eq:kineticstability}
 \sup_{u,v}d_2(z_v(u),z'_v(u))^2
 \le16B_I\left\{\frac{K(z)+K(z')}{2}-K(k)\right\},
 \qquad B_I=\int_I\sech u\,\dd u.
\end{equation}
This follows from \eqref{eq:polygonmidpoint} and
$\sum\Delta_j^2/W_j\le B_I$. It proves uniqueness of the relaxed
kinetic minimum.

Sample $k$ on refining partitions and apply $\Pi$ to its nodes.
The projected polygons have the same ordered endpoints and no larger
kinetic cost by Lemma~\ref{lem:pairretractions}. Their costs still
converge to $K(k)$. Inequality~\eqref{eq:kineticstability} forces
them to converge uniformly to $k$, and closedness of $\Dord$ makes
$k$ ordered. This argument does not identify a projected geodesic
with the geodesic joining its projected endpoints.

Finally, differentiating \eqref{eq:kineticbridge} gives
\[
 w_{vi}''=-\tanh u\,\sin(w_{vi}')\cos(w_{vi}'),
 \qquad |w_{vi}''|\le\tfrac12.
\]
The scalar interpolation error from the affine chord is at most
$|I|^2/16$ in each coordinate. Summing squares proves
\eqref{eq:bridgechord}.
\end{proof}

\begin{proposition}\label{prop:quadraticchord}
For $I=[l,r]$, let $\gamma$ be the affine endpoint chord of $c_*|I$.
With $M=\max_I R$ and $m=\min_I R$,
\begin{equation}\label{eq:quadraticchord}
 \sup_{u\in I}d_2((c_*)_v(u),\gamma_v(u))
 \le\sqrt n\left(16q\frac M m+\frac1{16}\right)|I|^2.
\end{equation}
In particular, for $|I|\le1$ the coefficient is bounded by
$C=\sqrt n(16qe+1/16)$, independently of the location of $I$ and
the boundary triple.
\end{proposition}

\begin{proof}
First work over $L$, using Proposition~\ref{prop:extensionminimum}.
The strict local margin makes the kinetic bridge $k$ feasible for
the endpoints of $c_*|I$. Put
$V(c)=q\int_I Rg(c)\,\dd u$ and
$\overline K(c)=\overline E_I(c)-V(c)$. This subtracts a continuous
potential from the relaxation; it does not presume a differential
action for $c_*$. Exact-endpoint recovery and
\eqref{eq:kineticstability} give, for
$D=\max_v\sup_I d_2((c_*)_v,k_v)$,
\[
 D^2\le8B_I\{\overline K(c_*)-K(k)\}.
\]
Local minimality and the orderedness of $k$ imply
\[
 0\le\overline K(c_*)-K(k)\le V(k)-V(c_*).
\]
The determinant logarithm is $\sqrt n$-Lipschitz for $d_2$.
Thus $|g(k)-g(c_*)|\le2\sqrt n D$, and
\[
 D^2\le16q\sqrt n B_I M|I|D.
\]
Divide by $D$ if it is nonzero and use $B_I\le |I|/m$.
Combining with \eqref{eq:bridgechord} proves
\eqref{eq:quadraticchord}. Since $|\dd(\log R)/\dd u|\le1$,
$M/m\le e^{|I|}$. Canonical extension is an isometry and preserves
the endpoint chord, so the result also holds for the original
$K$-norm curve.
\end{proof}

\subsection{Intrinsic spectral velocities}
For relative spectra we shall use the coordinatewise comparison
\begin{equation}\label{eq:coordinatecontinuity}
 \norm{\sigma(h,k)-\sigma(h',k')}_\infty
 \le d_\infty(h,h')+d_\infty(k,k').
\end{equation}
For split norms this follows by monotone comparison with scalar
multiples in the min--max description of each relative spectral
coordinate; see \cite[Lemma~2.23]{BE21}. Continuity extends it to all
norms. This estimate is stronger than the partial-sum bound needed
for earlier convexity arguments.

\begin{proposition}\label{prop:spectralvelocities}
At each vertex, the oriented relative secants of $c_*$ have a
common forward and backward decreasing limit $v(u)\in\RR^n$ at
every time. For $0<H\le1$,
\begin{equation}\label{eq:secantvelocity}
 \left\|\frac{\sigma(c_*(u),c_*(u+H))}{H}-v(u)\right\|_\infty
 \le4CH,
\end{equation}
with the same bound for the backward secant whose source is $u-H$
and target is $u$. The list $v$ is locally Lipschitz, with constant
$8C$. Its entries have a strict margin inside $(-\pi/2,\pi/2)$
on every compact interval. The determinant logarithm is $C^{1,1}$
locally and has derivative $\sum_i v_i$.
\end{proposition}

\begin{proof}
Fix a vertex and write $h(u)$ for its curve. Let
$f_u(H)=\sigma(h(u),h(u+H))/H$. The midpoint case of
\eqref{eq:quadraticchord}, spectral scaling on geodesics and
\eqref{eq:coordinatecontinuity} give
\[
 \norm{f_u(H/2)-f_u(H)}_\infty\le2CH.
\]
Summation on dyadic scales gives a limit and the error $4CH$.
For $H/2\le h\le H$, comparison with the chord at time $u+h$
gives $\norm{f_u(h)-f_u(H)}_\infty\le2CH$, identifying the
same limit along arbitrary scales. The backward secants have the
same construction. For $0<H\le1/2$, on $[u-H,u+H]$ the two half-chord spectra agree,
and the actual midpoint differs by at most $4CH^2$. Thus the
forward and backward secants differ by at most $8CH$, and their
limits coincide.

For $s<t$ the forward secant from $s$ to $t$ and the backward secant
at $t$ are identical. Applying their error estimates gives
$\norm{v(s)-v(t)}_\infty\le8C|s-t|$ for $t-s\le1$.
The strict margin follows from Proposition~\ref{prop:localmargin}.
Finally, the sum of the relative spectral coordinates is the
determinant increment. Its two difference quotients tend to
$\sum_i v_i(u)$ at every time, proving the last assertion.
\end{proof}

\begin{proposition}\label{prop:spectralaction}
On every compact interval the relaxed action of the actual minimum is
\begin{equation}\label{eq:spectralaction}
 \overline E_I(c_*|I)=\int_I R(u)
 \left\{\frac12\sum_{v=A,C}\sum_i\ell(v_{vi}(u))+qg(c_*(u))\right\}\dd u.
\end{equation}
The spectral positions $x_{vi}(u)=R(u)\tan v_{vi}(u)$ are locally
Lipschitz.
\end{proposition}

\begin{proof}
On a partition cell of length $\Delta$, its relative spectral
quotient differs by $O(\Delta)$ from the velocity list at the left
endpoint, and hence at every time in that cell. The strict compact
margin makes $\ell$ uniformly continuous on all these lists.
The partition lower-bound argument in
Lemma~\ref{lem:regularaction} therefore applies to any polygonal
sequence converging to $c_*$, without a fixed apartment for the
limit, and gives the lower bound in \eqref{eq:spectralaction}.

For recovery, sample $c_*$ and interpolate canonically. All-norm
nodes are permitted in the local relaxation. Order and the strict
metric speed bound make these finite-action polygons, with exact
endpoints. The same spectral-quotient estimate and the exact
segment weights give convergence to the displayed integral. The
potential converges uniformly as before. The final assertion follows
from the locally Lipschitz velocity lists and their strict margins.
\end{proof}

The velocities and positions in these propositions are numerical lists.
Their regularity alone does not identify embedded subspaces or make the
positions straight. Those are the remaining algebraic and variational
steps in the proof.
